\documentclass[12pt,a4paper]{report}

\usepackage[utf8]{inputenc}
\usepackage[T1]{fontenc}
\usepackage{geometry}
\usepackage{graphicx}
\usepackage{xcolor}
\usepackage{amsmath,amssymb}
\usepackage{booktabs}
\usepackage{float}
\usepackage{hyperref}
\usepackage[numbers,sort&compress]{natbib}

\geometry{top=2.5cm,bottom=2.5cm,left=2.5cm,right=2.5cm}

\title{Transient Grating Results Report}
\author{Manuel Fernando S\'anchez Alarc\'on}
\date{\today}

\begin{document}

\maketitle
\tableofcontents
\newpage

\chapter{Introduction}
This report summarizes the current status of transient grating analysis.
It is intended as the companion document for the code and notebook workflow in this repository.

\chapter{Data and Workflow}
\section{Dataset}
Describe here the scans, energies, and intensities used in the current analysis campaign.

\section{Processing Pipeline}
Document the preprocessing and fitting steps implemented in \texttt{src/tg\_analysis.py}.

\chapter{Results}
\section{Main Observations}
Add quantitative findings and their uncertainty estimates.

\section{Figures}
Insert figures exported from the analysis scripts and notebook.
For example:
\begin{figure}[H]
    \centering
    % \includegraphics[width=0.8\textwidth]{figures/example.png}
    \caption{Placeholder for transient grating analysis output.}
    \label{fig:tg-placeholder}
\end{figure}

\chapter{Conclusions}
Summarize the physical interpretation of the fitted parameters and open points for follow-up.

\bibliographystyle{unsrtnat}
\bibliography{references}

\end{document}
